
Vừa đọc trên Fb của GS Trần Văn Nhung và được biết GS Phan Đình Diệu đã mất. 

Tôi tốt nghiệp đại học toán Đại học Tổng hợp năm 1976 và được chọn ở lại trường làm giảng viên, nhưng được yêu cầu chuyển sang ngành Máy tính điện tử. Thời gian học máy tính ở khoa Toán, tôi chỉ biết một chút về lập trình và nguyên lý máy tính. 

Theo quy định thời đó, cán bộ giảng dạy phải tập sự hai năm và phải thi 4 chuyên đề. Ngành này khá mới mẻ, lúc đó ở bộ môn toán học tính toán chưa có ai là TS ngành máy tính (lúc đó gọi là chuyên ngành đảm bảo toán học cho máy tính) ngoài thầy Nguyễn Bá Hào thì lúc đó đã chuyển đi nơi khác . Chủ nhiệm Khoa (thày Hạp) mới gọi tôi lên đưa cho tôi một cái thư tay bảo: "Em sang bên Viện khoa học Tính toán và Điều khiển gặp thày Diệu, đây là thư tôi nhờ thày hướng dẫn cho em trong thời gian tập sự". 

Lúc bấy giờ Viện khoa học tính toán và điều khiển mới được thành lập, đóng đô ở Đồi Thông. Cứ sáng thứ 7 ở đó có xemina. Hầu hết những hành trang ban đầu về tin học tôi có là từ những xê mi na này. Thày Diệu hỏi han một chút và bảo: " Mình không có nhiều thì giờ lắm, Quốc đọc những cuốn sách này và thi với mình". Cuốn sách ông đưa là cuốn Ngôn ngữ hình thức và Chương trình dịch.   Hôm thi cũng thi trong lúc đang xê mina. Ông viết sẵn đề lên một nửa tờ A4 và tôi làm bài thi trong khi mọi người xê mi na. Tôi thi với thày Diệu 2 chuyên đề, còn hai chuyên đề khác thi với thày ở trong khoa.

Ông là người cực kỳ uyên bác và khúc triết. Nhưng bài ông trình bày trong xemina đều rất sâu sắc. Thời làm TS Khoa học bên Nga, thày làm về toán học kiến thiết, một thứ toán có liên quan mật thiết đến lý thuyết tính toán. Toán học kiến thiết không thừa nhận luật phản chứng, không thừa nhận vô hạn tiềm năng, chỉ thừa nhận vô hạn thực tai, không thừa nhận sự tồn tai theo kiểu kinh điển mà tồn tại tức là phải có một quy trình chỉ ra được sau hữu hạn bước. Điều này thật phù hợp với tư duy máy tính. Nghe nói, khi thành lập Viện Khoa học Tính toán và Điều khiển, người ta giao cho TSKH Nguyễn Thúc Loan làm Viện trưởng và chính ông Loan sau một thời gian ở vị trí này đã nói, xin để anh Diệu làm Viện trưởng, anh ấy làm sẽ tốt hơn tôi.

Thày xứng đáng được coi là người khai sinh ra ngành Tin học của Việt Nam. Thời đó tại Đồi Thông, VN là nước Đông Nam Á đầu tiên chế tạo ra máy vi tính theo đúng nghĩa. Rất nhiều các chuyên gia đầu ngành Tin học được đào tạo ở đây và ở Pháp qua đường hợp tác giữa VN và Pháp do Viện chủ trì.

Sau này do có sự bất đồng với lãnh đạo Viện hàn lâm KHVN, thày Diệu chuyển về làm cho Chương trình Quốc gia về CNTT giai đoạn 1996-2000 (thường được gọi là IT2000). Chương trình này đã đặt ra những nền tảng cho sự phát triển CNTT của VN sau này trong đó có việc thành lập các khoa CNTT trọng điểm, việc xây dựng các CSDL quốc gia, việc tin học hóa hành chính nhà nước. Năm kia tôi đọc lại các tài liệu về chương trình này và vẫn không hết ngạc nhiên về tầm nhìn của GS mặc dù đã qua 20 năm. 

Sau này ông xin về Đại học Quốc gia, Ông muốn tiếp tục cống hiến trong lĩnh vực giảng dạy. ĐHQG muốn ông về làm việc trực tiếp ở Khoa CNTT nhưng khoa tôi lại muốn ông sinh hoạt ở ĐHQGHN và chỉ giảng dạy ở Khoa, Khoa có một phòng làm việc dành riêng cho ông. Tôi có hỏi một lãnh đạo Khoa khi đó rằng tại sao lại đối xử như vậy thì được trả lời (chẳng hiểu có thật không) rằng "Tầm cỡ ông Diệu là phải có người giúp việc, có xe đưa xe đón thì khoa mình lo sao được", sau lại nói nhỏ  "phức tạp lắm đấy, mày có biết ai quan hệ với ông Diệu đều bị để ý không?"

Thời đó thày Diệu rất nổi tiếng theo kiểu một nhân vật bất đồng chính kiến khi ông từng nói về Lê Duẩn, "Ông Lê Duẩn là người rất vĩ đại, đã từng lãnh đạo Việt Nam đánh thắng Mỹ, thống nhất đất nước, nhưng sẽ vĩ đại hơn nếu từ chức TBT". Là đại biểu quốc hội, Ông có một lần đọc tham luận tại Quốc hội, một bản tham luận có tính phản biện cao, gây tiếng vang lớn mà thoát được sự kiểm duyệt trước. Ông luôn là người trăn trở với đất nước.

Sau khi bị đột quỵ, ông ít khi xuất hiện. Bức ảnh dưới đây là bức ảnh cuối cùng tôi chụp với ông trong buổi kỷ niệm 25 năm ngày thành lập Hội Tin học Việt Nam mà ông là Chủ tịch đầu tiên. Hồi đó ông đã yếu, đi lại chị Hương  phải dìu nhưng ông vẫn minh mẫn. 

Lần cùng với một số cán bộ khoa CNTT đến khoa cấp cứu A9 của bệnh viện Bạch Mai thăm ông, tôi phải đứng ngoài không nhìn thấy ông, nghe nói ông đang thiếp đi, tôi chỉ nói chuyện được một chút với con trai ông đang ở bệnh viện trông bố.

Xin chúc GS Phan Đình Diệu an nghỉ trong sự quý trọng và biết ơn của cộng đồng những người làm Công nghệ Thông tin. 
Vĩnh biệt giáo sư, 
Xin gửi lời chia buồn sâu sắc tới gia đình GS.